% r06_fill_DS: proof of Proposition rot:hyp:DS, the dead-shell chain in the exterior-shell geometry.
% Input from sec/r06_rot_shell2.tex right after the statement of \Cref{rot:hyp:DS}.

\begin{proof}[Proof of \Cref{rot:hyp:DS}]
Fix a two-level configuration with the exterior shell: the walled torus $T=D_{J'}$ of $\fc=(J',\psi,n)$, the source box $\Hboxp=\Qb{b'}$ with $b'\ge b+3$, the sharp profile $\chi=\mathbf 1_U$ with $U=T\setminus\Hboxp$, so that the dead set is $\Dchi=U$, and a King cube $C_L$ with $\Hboxp\subset C_L\subset T$, $\dist_\infty(\Hboxp,\R^3\setminus C_L)=D(\fs)+3$, $D(\fs)\ge\max(1,3/\lmu)$ and $\Nr1{\bar C_L}\subset(-25MJ',25MJ')^3$; let $f$ be a real $2$-form supported in $\Qb b$, $g=0$, and $k$ with $\eps=\eps_k\le\min(1/\lmu,1/4)$. Write $\Theta\defeq e^{\mathrm iF_{\Lk}(f)}$ and $E\defeq T\setminus C_L$ for the set of sites of the torus outside $C_L$. Since $\supp f\subset\Qb b$ and $b'\ge b+3$, every plaquette of $f$ lies in $\Hboxp$ at $\ell^\infty$-distance $>1$ from $\partial\Hboxp$. The dead-shell state $\Psik(n,\infty)=\langle\Theta\rangle$ is the walled torus state of $\fc$ with the matter killed on $U\cap\Lk$ (\Cref{rot:def:twolevel}): the single-site measure at every site of $U$ is the point mass at $\phi=0$, every hopping term with an endpoint in $U$ is reduced to its diagonal part at the other endpoint, and the photon is the periodic photon of $T$ with mass $\lmu$ (\Cref{rot:def:tori}(d)). The matter alive is that of $\Hboxp\cap\Lk$, with the wall $n\psi$ in its single-site measures; since $\psi=0$ on $\Nr1{\Qb b}$ and $\psi\ge1$ on $\Nr1U$, the wall is a fixed non-negative function on $\Hboxp$. The two differences from the cubic shell of \Cref{rot:def:shell} are that the dead region $U$ now extends to the whole exterior of $\Hboxp$ in $T$, so that no live matter sits outside $\Hboxp$, and that $\partial U=\partial\Hboxp$ is a single convex surface. Both only simplify the three steps.

\emph{Step 1 (pinning).} For $s\ge0$ add $s|A_\beta(y)|^2$ to the action for every photon variable $A_\beta(y)$ on a bond with an endpoint in $E$, that is for every $(\beta,y)$ with $y\in E$ or $y+\eps e_\beta\in E$. The resulting measures are BFS measures of King type (the paragraph after \Cref{rot:def:Zun}: adding $s|\ell(A)|^2$ to the action is one of the admitted operations), so by \Cref{lem:s02:BFS}(a) the expectation of $\Theta$, which is of class $\cI$, does not decrease in $s$; let $s\to\infty$. In the limit the pinned variables are fixed at $0$, and by the Gaussian structure of the photon part the law of the remaining variables $A_\alpha(x)$, $x,x+\eps e_\alpha\in C_L$, is the periodic photon of $T$ conditioned on the vanishing of all the others: this is the law ($\natural$) of \Cref{rot:lem:GB}, whose cell contains $\Nr1{C_L}$ by hypothesis. In the pinned measure the matter outside $\Hboxp$ is dead, so the hopping terms with both endpoints in $U$ contribute nothing (their cosine parts vanish at $\phi=0$ and their diagonal parts vanish as well), and a hopping term between a site of $\Hboxp$ and a dead neighbour in $U$ is its diagonal part $\frac\eps2r_x^2$ at the live endpoint, a Dirichlet term independent of $A$. In particular no hopping term couples $E$ to a live site, since every site of $C_L$ adjacent to $E$ lies in $C_L\setminus\Hboxp\subset U$ and is dead, and the matter on $E$ is a product of point masses, a factor $1$. What remains is exactly the $\natural$-measure of \Cref{rot:lem:GB}: the matter on $\Hboxp\cap\Lk$ with BFS single-site measures (the wall $n\psi$ and the Dirichlet terms toward $U$ included), hopping terms only between neighbouring sites of $\Hboxp$, coupled to the photon variables of the bonds of $\Hboxp$ as in King's action, and the photon law ($\natural$) on the bonds of $C_L$. Hence
\[
\Psik(n,\infty)\ \le\ \langle\Theta\rangle_\natural .
\]

\emph{Step 2 (the photon switch across the collar).} The collar $C_L\setminus\Hboxp\subset U$ carries no live matter, and $\dist_\infty(\Hboxp,\R^3\setminus C_L)=D(\fs)+3$ with $D(\fs)\ge\max(1,3/\lmu)$ and $\eps\le\min(1/\lmu,1/4)$. So \Cref{rot:lem:GB} applies to the pair of measures with the common matter part of Step~1 and the photon parts ($\natural$) and (N), the latter King's Neumann photon on $C_L\cap\Lk$; its hypotheses on the observable hold with $c=0$ and $g=0$, $f$ being supported on plaquettes of $\Hboxp$. It gives
\[
\langle\Theta\rangle_\natural\ \le\ \langle\Theta\rangle_N+\langle|\Theta|^2\rangle_N^{1/2}\big(e^{\tau}-1\big)^{1/2}
=\langle\Theta\rangle_N+\big(e^{\tau}-1\big)^{1/2},\qquad
\tau\le C_{\rm GB}\,\srw\,\vol_\eps(\Nr2{\Hboxp})\,D(\fs)^{-1}e^{-\frac32m_\eps D(\fs)},
\]
since $|\Theta|^2=1$ at $g=0$. With $m_\eps\ge0.96\lmu$ and $\vol_\eps(\Nr2{\Hboxp})\le\vol(\Nr3{\Hboxp})$ for $\eps\le1$ this is $\tau\le\tauGB(\fs)$, the constant \eqref{rot:eq:tau} computed from $\Hbox=\Hboxp$ and $D(\fs)$, as in Step~2 of the proof of \Cref{rot:thm:SHtoTIGHT}. The geometry of the shell outside $C_L$ plays no role here: the lemma sees only the matter on $\Hboxp$ and the photon on $C_L$.

\emph{Step 3 (to the comparison cube).} The $N$-measure of Step~2 is the Neumann--Dirichlet state on $X_L=C_L\cap\Lk$ of \Cref{rot:def:tori}(a) with the wall $n\psi$ in its single-site measures and the matter killed on $(C_L\setminus\Hboxp)\cap\Lk$: King's Neumann photon on $C_L$, Dirichlet terms for the matter of $\Hboxp$ toward its dead neighbours and, for the sites of $C_L$ adjacent to $\R^3\setminus C_L$, toward the outside (these sites are dead, so the two readings of their Dirichlet terms coincide). By \Cref{lem:s02:Kprime}, applied with the set $(C_L\setminus\Hboxp)\cap\Lk$ and the rate $1$ on it, this state is the limit $s\to\infty$ of the walled Neumann--Dirichlet states on $X_L$ in which the single-site measures on $(C_L\setminus\Hboxp)\cap\Lk$ are multiplied by $e^{-s\eps^3r_x^2}$; $|\Theta|\le1$ is of the form admitted there. By \Cref{lem:s02:BFS}(c) the expectation of $\Theta$ does not decrease when $s$ decreases to $0$ (the single-site weights $e^{-s\eps^3r_x^2}$ are removed), nor when the wall factors $e^{-\eps^3n\psi(x)r_x^2}$, $x\in X_L$, are then removed (the same clause, the wall being a non-negative multiplier of the form $e^{-sr^2}$ at each site). The result is the comparison cube $\ava\cdot$ on $X_L$ (\Cref{def:s02:acube}): the Neumann photon on $C_L$, the scalar alive on all of $X_L$ with the Dirichlet condition outside $C_L$, no wall. In it $\Theta$ is the same function of the shared variables, every plaquette of $f$ lying in $\Hboxp$ at $\ell^\infty$-distance $>1$ from $\R^3\setminus C_L$. Hence $\langle\Theta\rangle_N\le\ava{e^{\mathrm iF(f)}}|_{J=L}=Z_{k,L}(\mathrm if)$, and with Steps~1--2
\[
\Psik(n,\infty)\ \le\ Z_{k,L}(\mathrm if)+\big(e^{\tauGB(\fs)}-1\big)^{1/2},
\]
which is \eqref{rot:eq:deadbound} at $g=0$ for the exterior-shell configuration. The three steps use the shape of $U$ only through two facts, that every site of $C_L$ adjacent to $E$ is dead (Step~1) and that the collar between $\Hboxp$ and $\partial C_L$ carries no live matter (Step~2); both hold because $U\supset C_L\setminus\Hboxp$. The two-level sandwich \Cref{rot:lem:TL}, which \Cref{rot:thm:REStoTIGHT} applies to this configuration, needs $\chi\le\psi$ pointwise, that is $\psi\ge1$ on $U=T\setminus\Hboxp$; this is the plateau condition $\psi\ge1$ on $\Nr1U$ of \Cref{rot:def:shell}, kept for the exterior shell, and it is compatible with $\psi=0$ on $\Nr1{\Qb b}$ because $\dist_\infty(U,\Qb b)\ge3$.
\end{proof}

With \Cref{rot:hyp:DS} in place, the proof of \Cref{rot:thm:REStoTIGHT} runs verbatim for the exterior-shell configurations of \Cref{rot:thm:shell}: its only use of the shell geometry is the dead-shell bound \eqref{rot:eq:deadbound}, which the proposition supplies, and \Cref{rot:lem:TL}, which holds for the sharp profile by \Cref{rot:rem:Eproof}(ii).
